\scanpage{14}
\mspage{4}

\[
\text{a)}\quad \varphi_!(e_{\widehat A})\longrightarrow e_{\mathcal X}\quad\text{est épi}
\]
\[
\text{b)}\quad \varphi_!(\alpha\times\beta)\longrightarrow
\varphi_!(\alpha)\times\varphi_!(\beta)\quad\text{est épi}
\qquad \forall\alpha,\beta\text{ dans }\operatorname{Ob}A
\]
\[
\text{c)}\quad \varphi_!(\alpha\times_{\alpha_0}\beta)\longrightarrow
\varphi_!(\alpha)\times_{\varphi_!(\alpha_0)}\varphi_!(\beta)\quad\text{est épi}
\]
$\forall$ diagramme
\diagram{\begin{tikzcd}[ampersand replacement=\&,column sep=small,row sep=small]
\alpha\arrow[dr]\&\&\beta\arrow[dl]\\
\&\alpha_0\&
\end{tikzcd}}
dans $A$.

Un cas particulièrement agréable où ces conditions sont satisfaites, c'est quand $\varphi_!$ est exact à gauche, i.e. qu'il correspond \deleted{\unsure{à}} à un morphisme de topos
\[
f:\mathcal X\longrightarrow\widehat A
\]
par la formule
\[
\varphi_!=f^*.
\]
Alors dans a) b) c) les flèches épi en question sont même des iso, et je présume que l'inverse est vrai également. Oui -- et c'est essentiellement \unsure{evident|evid\textsuperscript{t}}, en écrivant canoniquement, pour un $F\in\operatorname{Ob}\widehat A$, $F=\varinjlim_{A/F}\alpha$ \unsure{$\ldots$}\unsure{.} On voit

\deleted{Mais \unsure{peut-on avoir}}
\inserted{de la $=$ façon pour} b) et c)

\margin{13}
